\documentclass[french]{article}
\usepackage[utf8]{inputenc}
\usepackage[T1]{fontenc}
\usepackage{lmodern}
\usepackage{geometry}
\usepackage{graphicx}
\usepackage{babel}
\usepackage{amsmath}
\usepackage{amsthm}
\usepackage{amssymb}
\usepackage{hyperref}
\usepackage{stmaryrd}
\usepackage{enumitem}
\usepackage{tcolorbox}
\usepackage{dsfont}
\usepackage{multirow}
\usepackage{etoc}
\usepackage{titlesec}
\usepackage{esint}
\usepackage{cancel}
\usepackage{mathdots}

\setlist[itemize]{label=$\bullet$}


\renewcommand{\P}{\mathbb{P}}
\newcommand{\N}{\mathbb{N}}
\newcommand{\Z}{\mathbb{Z}}
\newcommand{\Q}{\mathbb{Q}}
\newcommand{\R}{\mathbb{R}}
\newcommand{\C}{\mathbb{C}}
\newcommand{\K}{\mathbb{K}}
\renewcommand{\L}{\mathbb{L}}
\newcommand{\U}{\mathbb{U}}
\newcommand{\st}{^*}
\newcommand{\tq}{\middle|}
\newcommand{\inv}{^{-1}}
\newcommand{\E}{\mathbb{E}}
\newcommand{\F}{\mathbb{F}}
\newcommand{\V}{\mathbb{V}}
\renewcommand{\ker}{\text{Ker}}
\newcommand{\im}{\text{Im}}
\newcommand{\card}{\text{Card}}
\newcommand{\Tr}{\text{Tr}}

\title{TD Euclidiens\\ Exercice 17}
\date{}
\begin{document}
\maketitle

\section{Énoncé}
Soit $x-0$ un vecteur non nul d'un espace euclidien $E$ ainsi que $a$ et $b$ deux éléments de $E$. Donner une CNS pour qu'il existe $u\in\mathcal{L}(E)$ tel que $u(x_0)=a$ et $u\st(x_0)=b$.
\section{Solution}
Il est nécessaire que $(a\mid x_0)=(b\mid x_0)$ par définition de l'adjoint. Voyons en quoi cela est suffisant. Supposons dans un premier temps $\lVert x_0\rVert=1$. On considère $\mathcal{B}=(x_0,e_2,\ldots,e_n)$ une BON de $E$. Alors on regarde la matrice $A$ qui a pour première ligne les $(b\mid b_j)$, pour première colonne les $(a\mid b_i)$ (bien définie car $(a\mid x_0)=(b\mid x_0)$) et des $0$ partout ailleurs. On considère alors l'unique $u\in\mathcal{L}(E)$ tel que la matrice de $u$ dans $\mathcal{B}$ soit $A$, qui convient. Dans le cas général, on considère $u$ construit comme précédemment mais avec cette fois $$\left(\frac{a}{\lVert x_0\rVert}\mid\frac{x_0}{\lVert x_0\rVert}\right)=\left(\frac{b}{\lVert x_0\rVert}\mid\frac{x_0}{\lVert x_0\rVert}\right)$$ et alors par linéarité un tel $u$ convient.
\end{document}
